\section{Related Work}

\paragraph{Carbon-aware elasticity.}
CarbonScaler varies a batch workload's cloud allocation with carbon intensity
and uses a greedy marginal-resource allocation algorithm
~\cite{hanafy2023carbonscaler}. It already establishes resource elasticity as a
carbon-control mechanism. \name{} addresses the additional delay between
requesting resources and obtaining them, as well as repeated admission for
checkpointed chunks. A direct transplant of cloud autoscaling would change
the available controls; our controls retain repeated admissions and charge every queue wait. Studies of temporal and spatial shifting
also show that flexibility and overhead constrain attainable carbon savings
~\cite{sukprasert2024limitations}. We report the scope of the opportunity rather
than assume that lower-carbon timing is always available.

\paragraph{User-side provisioning.}
QBETS estimates wait bounds from historical time series~\cite{nurmi2007qbets};
our main policy learns completed-run costs without an admission estimator.
Mirage uses learning to provision follow-up resources for low-interruption
services on batch GPU clusters~\cite{ding2023mirage}. We use its trace-replay
infrastructure, but optimize the node count of a ready training chunk under a
completion budget. Scheduler-level methods such as DRAS act on the scheduling
policy itself~\cite{fan2021deep}. Our policy has no such authority and cannot
observe hidden priority or reservation state.

\paragraph{Constrained learning and accounting uncertainty.}
Constrained policy optimization separates a resource objective from a service
constraint~\cite{achiam2017cpo}; PPO supplies a practical actor--critic optimizer
~\cite{schulman2017ppo}. We use these established ideas, without claiming a new
general-purpose RL optimizer or a constraint-satisfaction theorem. The method
contribution is the request-conditioned decision representation and its
coupling to a tractable power-scenario objective. Unlike operational-carbon
measurement, the latter assesses consequences conditional on a model. An
algebraic sensitivity certificate does not validate the model's physical
parameters.
